\chapter{2004年安徽卷（春理）}

\section{单选题}

\begin{problem}
\(\displaystyle\frac{5(4+\i)^2}{\i(2+\i)}=\)（\quad）

\choices
    {\(5(1-38\i)\)}
    {\(5(1+38\i)\)}
    {\(1+38\i\)}
    {\(1-38\i\)}
\end{problem}

\begin{answer}
D
\end{answer}

\begin{solution}
先计算分子和分母，得 \((4+\i)^2=15+8\i\)，且 \(\i(2+\i)=-1+2\i\). 因此
\[
\frac{5(4+\i)^2}{\i(2+\i)}
=\frac{75+40\i}{-1+2\i}
=\frac{(75+40\i)(-1-2\i)}{(-1+2\i)(-1-2\i)}
=\frac{5-190\i}{5}=1-38\i.
\]
故选 D.
\end{solution}

\begin{problem}
不等式 \(\left|2x^2-1\right|\leqslant1\) 的解集为（\quad）

\choices
    {\(\{x\mid-1\leqslant x\leqslant1\}\)}
    {\(\{x\mid-2\leqslant x\leqslant2\}\)}
    {\(\{x\mid0\leqslant x\leqslant2\}\)}
    {\(\{x\mid-2\leqslant x\leqslant0\}\)}
\end{problem}

\begin{answer}
A
\end{answer}

\begin{solution}
由绝对值不等式得 \(-1\leqslant2x^2-1\leqslant1\). 两边分别加 \(1\)，再除以正数 \(2\)，得 \(0\leqslant x^2\leqslant1\). 因而 \(-1\leqslant x\leqslant1\)，故选 A.
\end{solution}

\begin{problem}
已知 \(F_1\)，\(F_2\) 为椭圆 \(\displaystyle\frac{x^2}{a^2}+\frac{y^2}{b^2}=1\)（\(a>b>0\)）的焦点，\(M\) 为椭圆上一点，\(MF_1\) 垂直于 \(x\) 轴，且 \(\angle F_1MF_2=60^\circ\)，则椭圆的离心率为（\quad）

\choices
    {\(\dfrac12\)}
    {\(\dfrac{\sqrt2}{2}\)}
    {\(\dfrac{\sqrt3}{3}\)}
    {\(\dfrac{\sqrt3}{2}\)}
\end{problem}

\begin{answer}
C
\end{answer}

\begin{solution}
设椭圆的半焦距为 \(c\)，则 \(F_1F_2=2c\). 因为 \(F_1F_2\) 在 \(x\) 轴上，而 \(MF_1\perp x\) 轴，所以三角形 \(F_1MF_2\) 在 \(F_1\) 处为直角. 又 \(\angle F_1MF_2=60^\circ\)，设 \(MF_1=u\)，则
\[
MF_2=2u,
\qquad
F_1F_2=\sqrt{(2u)^2-u^2}=\sqrt3u.
\]
由椭圆的焦半径和为定值，得 \(MF_1+MF_2=2a\)，所以 \(3u=2a\). 另一方面 \(2c=\sqrt3u\)，从而 \(2c=\dfrac{2\sqrt3}{3}a\)，即
\[
e=\frac ca=\frac{\sqrt3}{3}.
\]
故选 C.
\end{solution}

\begin{problem}
\(\displaystyle\lim_{n\to\infty}\frac{(n-2)^2(2+3n)^3}{(1-n)^5}=\)（\quad）

\choices
    {\(0\)}
    {\(32\)}
    {\(-27\)}
    {\(27\)}
\end{problem}

\begin{answer}
C
\end{answer}

\begin{solution}
分子分母同除以 \(n^5\)，得
\[
\frac{(n-2)^2(2+3n)^3}{(1-n)^5}
=\frac{\left(1-\frac2n\right)^2\left(3+\frac2n\right)^3}{\left(\frac1n-1\right)^5}.
\]
当 \(n\to\infty\) 时，分子趋于 \(1\cdot3^3=27\)，分母趋于 \((-1)^5=-1\)，所以极限为 \(-27\)，故选 C.
\end{solution}

\begin{problem}
等边三角形 \(ABC\) 的边长为 \(4\)，\(M\)，\(N\) 分别为 \(AB\)，\(AC\) 的中点，沿 \(MN\) 将 \(\triangle AMN\) 折起，使得面 \(AMN\) 与面 \(MNCB\) 所成的二面角为 \(30^\circ\)，则四棱锥 \(A-MNCB\) 的体积为（\quad）

\choices
    {\(\dfrac32\)}
    {\(\dfrac{\sqrt3}{2}\)}
    {\(\sqrt3\)}
    {\(3\)}
\end{problem}

\begin{answer}
A
\end{answer}

\begin{solution}
因为 \(M\)，\(N\) 分别是 \(AB\)，\(AC\) 的中点，所以 \(MN\parallel BC\)，且 \(MN=2\). 三角形 \(AMN\) 为边长为 \(2\) 的等边三角形，故点 \(A\) 到直线 \(MN\) 的距离为 \(\sqrt3\). 四边形 \(MNCB\) 是梯形，其两底为 \(MN=2\)，\(BC=4\)，高为正三角形 \(ABC\) 高的一半，即 \(\sqrt3\)，所以
\[
S_{MNCB}=\frac{2+4}{2}\cdot\sqrt3=3\sqrt3.
\]
折叠后，点 \(A\) 到平面 \(MNCB\) 的距离等于点 \(A\) 到折痕 \(MN\) 的距离乘二面角的正弦，因此四棱锥的高为 \(\sqrt3\sin30^\circ=\dfrac{\sqrt3}{2}\). 所以
\[
V=\frac13S_{MNCB}\cdot\frac{\sqrt3}{2}=\frac32.
\]
故选 A.
\end{solution}

\begin{problem}
已知数列 \(\{a_n\}\) 满足 \(a_0=1\)，\(a_n=a_0+a_1+\cdots+a_{n-1}\)（\(n\geqslant1\)），则当 \(n\geqslant1\) 时，\(a_n=\)（\quad）

\choices
    {\(2^n\)}
    {\(\dfrac{n(n+1)}2\)}
    {\(2^{n-1}\)}
    {\(2^n-1\)}
\end{problem}

\begin{answer}
C
\end{answer}

\begin{solution}
当 \(n\geqslant2\) 时，分别利用相邻两项的定义，有
\[
a_n-a_{n-1}=a_{n-1},
\]
所以 \(a_n=2a_{n-1}\). 又由已知 \(a_1=a_0=1\)，因此数列从 \(a_1\) 开始是首项为 \(1\)、公比为 \(2\) 的等比数列，故
\[
a_n=2^{n-1}\quad(n\geqslant1).
\]
故选 C.
\end{solution}

\begin{problem}
若二面角 \(\alpha-l-\beta\) 为 \(120^\circ\)，直线 \(m\perp\alpha\)，则 \(\beta\) 所在平面内的直线与 \(m\) 所成角的取值范围是（\quad）

\choices
    {\((0^\circ,90^\circ]\)}
    {\([30^\circ,60^\circ]\)}
    {\([60^\circ,90^\circ]\)}
    {\([30^\circ,90^\circ]\)}
\end{problem}

\begin{answer}
D
\end{answer}

\begin{solution}
二面角的两个补角相等，所以平面 \(\alpha\)，\(\beta\) 的锐二面角为 \(60^\circ\). 直线 \(m\) 垂直于平面 \(\alpha\)，因此它与平面 \(\beta\) 所成的角为 \(90^\circ-60^\circ=30^\circ\). 在平面 \(\beta\) 内，取其内的直线方向逐渐转动，直线与 \(m\) 所成的角的最小值为 \(30^\circ\)，当该直线垂直于 \(m\) 在平面 \(\beta\) 内的射影时取最大值 \(90^\circ\). 故取值范围为 \([30^\circ,90^\circ]\)，选 D.
\end{solution}

\begin{problem}
若 \(f(\sin x)=2-\cos2x\)，则 \(f(\cos x)=\)（\quad）

\choices
    {\(2-\sin2x\)}
    {\(2+\sin2x\)}
    {\(2-\cos2x\)}
    {\(2+\cos2x\)}
\end{problem}

\begin{answer}
D
\end{answer}

\begin{solution}
在已知关系中令 \(x\) 换成 \(\dfrac\pi2-x\)，则 \(\sin\left(\dfrac\pi2-x\right)=\cos x\)，从而
\[
f(\cos x)=2-\cos\left(\pi-2x\right)=2+\cos2x.
\]
故选 D.
\end{solution}

\begin{problem}
直角坐标 \(xOy\) 平面上，平行直线 \(x=n\)（\(n=0\)，\(1\)，\(2\)，\(\ldots\)，\(5\)）与平行直线 \(y=n\)（\(n=0\)，\(1\)，\(2\)，\(\ldots\)，\(5\)）组成的图形中，矩形共有（\quad）

\choices
    {\(25\) 个}
    {\(36\) 个}
    {\(100\) 个}
    {\(225\) 个}
\end{problem}

\begin{answer}
D
\end{answer}

\begin{solution}
六条竖直直线中任取两条确定矩形的两条竖边，有 \(\mathrm C_6^2\) 种取法；六条水平直线中任取两条确定矩形的两条横边，也有 \(\mathrm C_6^2\) 种取法. 所以矩形总数为
\[
\mathrm C_6^2\mathrm C_6^2=15\times15=225.
\]
故选 D.
\end{solution}

\begin{problem}
已知直线 \(l:x-y-1=0\)，\(l_1:2x-y-2=0\). 若直线 \(l_2\) 与 \(l_1\) 关于 \(l\) 对称，则 \(l_2\) 的方程是（\quad）

\choices
    {\(x-2y+1=0\)}
    {\(x-2y-1=0\)}
    {\(x+y-1=0\)}
    {\(x+2y-1=0\)}
\end{problem}

\begin{answer}
B
\end{answer}

\begin{solution}
两直线 \(l\)，\(l_1\) 的交点为 \(P(1,0)\). 令平移后的坐标为 \(X=x-1\)，\(Y=y\)，则对称轴 \(l\) 的方程为 \(Y=X\)，直线 \(l_1\) 的方程为 \(Y=2X\). 关于 \(Y=X\) 对称就是交换 \(X\)，\(Y\)，所以直线 \(Y=2X\) 的对称直线满足 \(Y=\dfrac12X\). 化回原坐标，得 \(y=\dfrac12(x-1)\)，即
\[
x-2y-1=0.
\]
故选 B.
\end{solution}

\begin{problem}
已知向量集合 \(M=\{\bs{a}\mid \bs{a}=(1,2)+\lambda(3,4),\lambda\in\R\}\)，\(N=\{\bs{a}\mid\bs{a}=(-2,-2)+\lambda(4,5),\lambda\in\R\}\)，则 \(M\cap N=\)（\quad）

\choices
    {\(\{(1,1)\}\)}
    {\(\{(1,1),(-2,-2)\}\)}
    {\(\{(-2,-2)\}\)}
    {\(\varnothing\)}
\end{problem}

\begin{answer}
C
\end{answer}

\begin{solution}
若向量属于交集，设对应参数分别为 \(\lambda\)，\(\mu\)，则其坐标满足
\[
\begin{cases}
1+3\lambda=-2+4\mu,\\
2+4\lambda=-2+5\mu.
\end{cases}
\]
即 \(3\lambda-4\mu=-3\)，\(4\lambda-5\mu=-4\). 第一式乘以 \(4\)，第二式乘以 \(3\) 后相减，得 \(\mu=0\)，再代回得 \(\lambda=-1\). 交点为 \((-2,-2)\)，所以
\[
M\cap N=\{(-2,-2)\}.
\]
故选 C.
\end{solution}

\begin{problem}
函数 \(y=\sin^4x+\cos^4x\) 的最小正周期为（\quad）

\choices
    {\(\dfrac\pi4\)}
    {\(\dfrac\pi2\)}
    {\(\pi\)}
    {\(2\pi\)}
\end{problem}

\begin{answer}
B
\end{answer}

\begin{solution}
利用平方和公式，得
\[
\sin^4x+\cos^4x
=(\sin^2x+\cos^2x)^2-2\sin^2x\cos^2x
=1-\frac12\sin^22x
=\frac34+\frac14\cos4x.
\]
因为 \(\cos4x\) 的最小正周期为 \(\dfrac\pi2\)，且系数不为零，所以原函数的最小正周期为 \(\dfrac\pi2\)，故选 B.
\end{solution}

\section{填空题}

\begin{problem}
抛物线 \(y^2=6x\) 的准线方程为\fillinblank{}.
\end{problem}

\begin{answer}
\(x=-3\)
\end{answer}

\begin{solution}
抛物线的标准方程为 \(y^2=4px\)，与 \(y^2=6x\) 比较得 \(p=\dfrac32\). 因而其准线方程为 \(x=-p\)，即 \(x=-3\).
\end{solution}

\begin{problem}
在 \(5\) 名学生（\(3\) 名男生，\(2\) 名女生）中安排 \(2\) 名学生值日，其中至少有 \(1\) 名女生的概率是\fillinblank{}.
\end{problem}

\begin{answer}
\(\dfrac7{10}\)
\end{answer}

\begin{solution}
从 \(5\) 名学生中任取 \(2\) 名，共有 \(\mathrm C_5^2=10\) 种等可能结果. 至少有 \(1\) 名女生包括一名女生一名男生或两名女生，共有
\[
\mathrm C_2^1\mathrm C_3^1+\mathrm C_2^2=6+1=7
\]
种. 因而所求概率为 \(\dfrac7{10}\).
\end{solution}

\begin{problem}
函数 \(y=\sqrt{x}-x\)（\(x\geqslant0\)）的最大值为\fillinblank{}.
\end{problem}

\begin{answer}
\(\dfrac14\)
\end{answer}

\begin{solution}
令 \(t=\sqrt{x}\)，则 \(t\geqslant0\)，且 \(x=t^2\). 函数值为
\[
y=t-t^2=\frac14-\left(t-\frac12\right)^2\leqslant\frac14.
\]
当 \(t=\dfrac12\)，即 \(x=\dfrac14\) 时等号成立，所以最大值为 \(\dfrac14\).
\end{solution}

\begin{problem}
若 \(\left(x+\dfrac1x-2\right)^n\) 的展开式中常数项为 \(-20\)，则自然数 \(n=\)\fillinblank{}.
\end{problem}

\begin{answer}
\(3\)
\end{answer}

\begin{solution}
先因式分解，得
\[
x+\frac1x-2=\frac{(x-1)^2}{x}.
\]
所以
\[
\left(x+\frac1x-2\right)^n=x^{-n}(x-1)^{2n}.
\]
要取常数项，必须从 \((x-1)^{2n}\) 中取出 \(x^n\) 项，其系数为 \((-1)^n\mathrm C_{2n}^n\). 因而
\[
(-1)^n\mathrm C_{2n}^n=-20.
\]
这要求 \(n\) 为奇数且 \(\mathrm C_{2n}^n=20\). 当 \(n=0\) 时常数项为 \(1\)，不合题意. 对 \(n\geqslant1\)，有

\[
\frac{\mathrm C_{2n+2}^{n+1}}{\mathrm C_{2n}^n}
=\frac{(2n+2)(2n+1)}{(n+1)^2}
=2\frac{2n+1}{n+1}>1,
\]

所以 \(\mathrm C_{2n}^n\) 随 \(n\) 严格递增. 又 \(\mathrm C_2^1=2\)，\(\mathrm C_6^3=20\)，故只有 \(n=3\) 能满足 \(\mathrm C_{2n}^n=20\)，唯一解为 \(n=3\).
\end{solution}

\section{解答题}

\begin{problem}
解关于 \(x\) 的不等式：\(\log_a^3x<3\log_ax\)（\(a>0\)，\(a\ne1\)）.
\end{problem}

\begin{solution}
由对数的定义域知 \(x>0\). 令 \(t=\log_ax\)，则当 \(x\) 遍历正数时，\(t\) 遍历全体实数. 原不等式化为
\[
t^3<3t\Longleftrightarrow t(t^2-3)<0\Longleftrightarrow t(t-\sqrt3)(t+\sqrt3)<0.
\]
在三个根 \(-\sqrt3\)，\(0\)，\(\sqrt3\) 处分区间讨论符号，得
\[
t\in(-\infty,-\sqrt3)\cup(0,\sqrt3).
\]
当 \(a>1\) 时，\(\log_a x\) 单调递增，故
\[
x\in(0,a^{-\sqrt3})\cup(1,a^{\sqrt3}).
\]
当 \(0<a<1\) 时，\(\log_a x\) 单调递减，故
\[
x\in(a^{\sqrt3},1)\cup(a^{-\sqrt3},+\infty).
\]
\end{solution}

\begin{problem}
已知正项数列 \(\{b_n\}\) 的前 \(n\) 项和 \(B_n=\dfrac14(b_n+1)^2\)，求 \(\{b_n\}\) 的通项公式.
\end{problem}

\begin{solution}
因为 \(B_1=b_1\)，所以
\[
b_1=\frac14(b_1+1)^2\Longleftrightarrow(b_1-1)^2=0.
\]
故 \(b_1=1\). 对 \(n\geqslant2\)，由 \(b_n=B_n-B_{n-1}\)，得
\[
4b_n=(b_n+1)^2-(b_{n-1}+1)^2.
\]
令 \(c_n=b_n+1\)，则 \(c_n>1\)，且上式化为
\[
c_{n-1}^2=c_n^2-4c_n+4=(c_n-2)^2.
\]
因此 \(c_{n-1}=|c_n-2|\)，即 \(c_n=2+c_{n-1}\) 或 \(c_n=2-c_{n-1}\). 因为 \(c_{n-1}>1\)，第二种情形给出 \(c_n<1\)，与 \(c_n>1\) 矛盾，所以只能有 \(c_n=2+c_{n-1}\). 又 \(c_1=b_1+1=2\)，归纳得 \(c_n=2n\). 因而
\[
b_n=c_n-1=2n-1.
\]
\end{solution}

\begin{problem}
已知 \(k>0\)，直线 \(l_1:y=kx\)，\(l_2:y=-kx\).
\begin{enumerate}
\item 证明：到 \(l_1\)，\(l_2\) 的距离的平方和为定值 \(a\)（\(a>0\)）的点的轨迹是圆或椭圆；
\item 求到 \(l_1\)，\(l_2\) 的距离之和为定值 \(c\)（\(c>0\)）的点的轨迹.
\end{enumerate}
\end{problem}

\begin{solution}
设动点为 \(P(x,y)\). 点 \(P\) 到两直线的距离分别为
\[
d_1=\frac{|kx-y|}{\sqrt{1+k^2}},
\qquad
d_2=\frac{|kx+y|}{\sqrt{1+k^2}}.
\]
\begin{enumerate}
\item 由平方和公式，得
\[
d_1^2+d_2^2
=\frac{(kx-y)^2+(kx+y)^2}{1+k^2}
=\frac{2(k^2x^2+y^2)}{1+k^2}=a.
\]
即
\[
\frac{x^2}{\dfrac{a(1+k^2)}{2k^2}}+\frac{y^2}{\dfrac{a(1+k^2)}2}=1.
\]
这是以原点为中心的椭圆；当两半轴相等，即 \(k=1\) 时，它是圆.

\item 利用恒等式 \(|u-v|+|u+v|=2\max\{|u|,|v|\}\)，取 \(u=kx\)，\(v=y\)，有
\[
d_1+d_2=\frac{2\max\{k|x|,|y|\}}{\sqrt{1+k^2}}=c.
\]
令 \(d=\dfrac{c\sqrt{1+k^2}}2\)，则轨迹满足
\[
\max\{k|x|,|y|\}=d.
\]
因此轨迹是矩形 \(|x|\leqslant\dfrac dk\)，\(|y|\leqslant d\) 的边界，即
\[
\begin{cases}
|x|=\dfrac dk,\enspace |y|\leqslant d,\\
|x|\leqslant\dfrac dk,\enspace |y|=d.
\end{cases}
\]
\end{enumerate}
\end{solution}

\begin{problem}
如图，已知三棱柱 \(ABC-A_1B_1C_1\) 中，底面边长和侧棱长均为 \(a\)，侧面 \(A_1ACC_1\perp\) 底面 \(ABC\)，\(A_1B=\dfrac{\sqrt6}{2}a\).

\begin{enumerate}
\item 求异面直线 \(AC\) 与 \(BC_1\) 所成角的余弦值；
\item 求证：\(A_1B\perp\) 面 \(AB_1C\).
\end{enumerate}

\bitmapfigure[max width=0.45\linewidth]{img/2004/anhui_science_spring/q20_fig1.png}
\end{problem}

\begin{solution}
建立空间直角坐标系，使底面 \(ABC\) 在 \(z=0\) 平面内，取 \(A(0,0,0)\)，\(C(a,0,0)\)，\(B\left(\dfrac a2,\dfrac{\sqrt3}{2}a,0\right)\). 由于底面为正三角形，且侧面 \(A_1ACC_1\) 垂直于底面并以 \(AC\) 为交线，可令 \(A_1=(u,0,v)\). 由侧棱长为 \(a\)，有
\[
u^2+v^2=a^2.
\]
又由 \(A_1B=\dfrac{\sqrt6}{2}a\)，得
\[
\left(u-\frac a2\right)^2+\frac34a^2+v^2=\frac32a^2.
\]
结合 \(u^2+v^2=a^2\)，化简得 \(u=\dfrac a2\)，进而 \(v^2=\dfrac34a^2\). 取 \(v=\dfrac{\sqrt3}{2}a\)，则
\[
A_1\left(\frac a2,0,\frac{\sqrt3}{2}a\right),
\quad
B_1\left(a,\frac{\sqrt3}{2}a,\frac{\sqrt3}{2}a\right),
\quad
C_1\left(\frac{3a}{2},0,\frac{\sqrt3}{2}a\right).
\]
\begin{enumerate}
\item 取 \(AC\) 的方向向量 \(\overrightarrow{AC}=(a,0,0)\)，取 \(BC_1\) 的方向向量
\[
\overrightarrow{BC_1}=\left(a,-\frac{\sqrt3}{2}a,\frac{\sqrt3}{2}a\right).
\]
于是
\[
\overrightarrow{AC}\cdot\overrightarrow{BC_1}=a^2,
\qquad
|\overrightarrow{AC}|=a,
\qquad
|\overrightarrow{BC_1}|=a\sqrt{\frac52}.
\]
所以异面直线所成角的余弦值为
\[
\frac{|\overrightarrow{AC}\cdot\overrightarrow{BC_1}|}{|\overrightarrow{AC}|\,|\overrightarrow{BC_1}|}=\sqrt{\frac25}.
\]

\item 取 \(A_1B\) 的方向向量 \(\overrightarrow{A_1B}=\left(0,\dfrac{\sqrt3}{2}a,-\dfrac{\sqrt3}{2}a\right)\). 平面 \(AB_1C\) 内有两条相交直线 \(AC\)，\(AB_1\)，且
\[
\overrightarrow{A_1B}\cdot\overrightarrow{AC}=0,
\qquad
\overrightarrow{A_1B}\cdot\overrightarrow{AB_1}=0.
\]
此外，\(A_1B\) 的中点
\[
Q\left(\frac a2,\frac{\sqrt3}{4}a,\frac{\sqrt3}{4}a\right)
\]
也是 \(AB_1\) 的中点，故 \(Q\) 在面 \(AB_1C\) 内. 因此 \(A_1B\) 同时垂直于平面 \(AB_1C\) 内两条相交直线 \(AC\)，\(AB_1\)，且与该平面相交，故 \(A_1B\perp\) 面 \(AB_1C\).
\end{enumerate}
\end{solution}

\begin{problem}
已知盒中有 \(10\) 个灯泡，其中 \(8\) 个正品，\(2\) 个次品. 现需要从中取出 \(2\) 个正品，每次取出 \(1\) 个，取出后不放回，直到取出 \(2\) 个正品为止. 设 \(\xi\) 为取出的次数，求 \(\xi\) 的分布列及 \(E(\xi)\).
\end{problem}

\begin{solution}
\(\xi\) 至少为 \(2\)，至多为 \(4\). 当前两次取出的都是正品时，停止于第 \(2\) 次，因此
\[
P(\xi=2)=\frac8{10}\cdot\frac79=\frac{28}{45}.
\]
停止于第 \(3\) 次，等价于前两次恰有一个正品、一个次品，且第三次为正品，所以
\[
P(\xi=3)=2\cdot\frac8{10}\cdot\frac29\cdot\frac78=\frac{14}{45}.
\]
停止于第 \(4\) 次，等价于前三次含有两个次品和一个正品，第四次必为正品，因此
\[
P(\xi=4)=3\cdot\frac8{10}\cdot\frac29\cdot\frac18=\frac1{15}=\frac3{45}.
\]
所以分布列可写为
\[
P(\xi=j)=
\begin{cases}
\dfrac{28}{45},&\enspace j=2,\\
\dfrac{14}{45},&\enspace j=3,\\
\dfrac3{45},&\enspace j=4,\\
0,&\enspace \text{其他}.
\end{cases}
\]
并且
\[
E(\xi)=2\cdot\frac{28}{45}+3\cdot\frac{14}{45}+4\cdot\frac3{45}=\frac{22}{9}.
\]
\end{solution}

\begin{problem}
已知抛物线 \(C:y=x^2+4x+\dfrac27\)，过 \(C\) 上一点 \(M\)，且与 \(M\) 处的切线垂直的直线称为 \(C\) 在点 \(M\) 的法线.
\begin{enumerate}
\item 若 \(C\) 在点 \(M\) 的法线的斜率为 \(-\dfrac12\)，求点 \(M\) 的坐标 \((x_0,y_0)\)；
\item 设 \(P(-2,a)\) 为 \(C\) 对称轴上的一点，在 \(C\) 上是否存在点，使得 \(C\) 在该点的法线通过点 \(P\)？若有，求出这些点，以及 \(C\) 在这些点的法线方程；若没有，请说明理由.
\end{enumerate}
\end{problem}

\begin{solution}
将抛物线方程写成
\[
y=(x+2)^2-\frac{26}{7},
\]
所以其对称轴为 \(x=-2\). 在点 \(M(x_0,y_0)\) 处，切线斜率为 \(2x_0+4=2(x_0+2)\).
\begin{enumerate}
\item 当法线斜率为 \(-\dfrac12\) 时，切线斜率为 \(2\)，所以
\[
2(x_0+2)=2,
\qquad x_0=-1.
\]
代入抛物线方程，得
\[
y_0=(-1)^2+4(-1)+\frac27=-\frac{19}{7}.
\]
故 \(M\left(-1,-\dfrac{19}{7}\right)\).

\item 先考虑顶点 \(M_0\). 当 \(x_0=-2\) 时，切线水平，法线为竖直直线 \(x=-2\)，它经过轴上任意点 \(P(-2,a)\). 因此顶点
\[
M_0\left(-2,-\frac{26}{7}\right)
\]
始终是满足条件的点，其法线方程为 \(x=-2\).

再设 \(x_0\ne-2\). 法线斜率为 \(-\dfrac1{2(x_0+2)}\)，法线过点 \(P(-2,a)\) 等价于
\[
a-y_0=-\frac1{2(x_0+2)}(-2-x_0)=\frac12.
\]
又 \(y_0=(x_0+2)^2-\dfrac{26}{7}\)，所以
\[
(x_0+2)^2=a+\frac{45}{14}.
\]
令 \(D=a+\dfrac{45}{14}\). 当 \(D>0\) 时，除顶点外还有两个点
\[
M_{\pm}\left(-2\pm\sqrt D,\,a-\frac12\right),
\]
它们处的法线分别为
\[
y-a=-\frac{x+2}{2\sqrt D},
\qquad
y-a=\frac{x+2}{2\sqrt D}.
\]
当 \(D\leqslant0\) 时，不存在除顶点以外的点. 综上，任意 \(a\) 都至少有顶点 \(M_0\) 满足条件；当 \(a+\dfrac{45}{14}>0\) 时，再增加上述两个点.
\end{enumerate}
\end{solution}
